% \begin{landscape}
    
    \begin{table*}[t]
        \begin{center}
        {\small
            \begin{tabular}{l|l} 
                \textbf{Problem}                                                            &   \textbf{Output Format}\\
                \hline
                \textbf{[1]} Who is in the same group as \texttt{(x)}?                                   & List of Vertex IDs    \\          
                \textbf{[2]} Is \texttt{(x)} likely to be know \texttt{(y)}?                             & Tuple of (Y/N, Confidence)  \\               
                \textbf{[3]} Count the vertices in the largest and                                       & Tuple of (small\_size,big\_size,[small],[big])\\
                ~~~~~~smallest disjoint segments.                                                  & ~ \\
                \textbf{[4]} \texttt{(x)}, \texttt{(y)} and \texttt{(z)} are in communities              & List of Vertex IDs\\
                ~~~~~~\texttt{(a)}, \texttt{(b)} and \texttt{(c)}. Who else is in \texttt{(c)}?   &~ \\
                \textbf{[5]} Who shares a segment with both \texttt{(a)} with                & List of Vertex IDs\\
                ~~~~~~property \texttt{(1)}and \texttt{(b)} with property \texttt{(2)}?          &\\
                \textbf{[6]} Which segment has the most vertices with                          & Tuple of (Group\_name, [Vertex IDs])\\
                ~~~~~~property \texttt{(3)}?                                                      & \\
                \textbf{[7]} What are the segments and who is in each?                               & List of Lists of Vertex IDs   \\ 
                \textbf{[8]} There are \texttt{(n)} segments, which segment                      & List of Lists of Vertex IDs\\
                ~~~~~~does each vertex belong to?                                                 &~\\                

            \end{tabular}
        \normalsize
        \caption{Example Problems for Population Segmentation} 
        \label{table:cd-problem-definition}}
        \end{center}
    \end{table*}
% \end{landscape}


% \begin{enumerate}
%     \item %Requires them to identify a `group' and check for membership by both. Could use CD, could do BFS to some threshold, and could just make a guess. 
%     \item  Return a tuple of the form (Yes/No, Confidence)? %Requires them figure out if they are affiliated. Might do SSSP, might do CD and see if in the same group
%     \item Return a tuple with two ints and two lists of form (Size\_big, Size\_small, [Big\_vertices], [Small\_vertices]) %Requires communities to be detected and ranked by size. 
%     \nrscomment{format the variables in a font with equal kearning}
%     \item  Return a list of Vertices. %Gives extra information, they may choose to use it or not.
%     \item Return a tuple of the form ([list\_of\_vertices], property).
%     \item What are the groups in the graph? Return a list of lists of vertices. % Execute community detection and return the result. 
%     \item There are (n) communities in this graph, named from largest to smallest $C_{1} ... C_{n}$. Which vertices belong to each group? Return a list of lists of vertices. % Execute CD with known number of communities. 
% \end{enumerate}

% \end{landscape}